% ============================================================
\chapter{Resumen}

Este trabajo presenta \textit{kubescan}, un sistema de detección automática de
cadenas de ataque en infraestructuras \textit{Kubernetes} que combina un clasificador
\textit{Random Forest} (Capa~1) con una \textit{Graph Attention Network}
(\acs{GAT}) (Capa~2) y optimización de pesos mediante Algoritmo Genético (Capa~3).

A partir de 2.827 manifiestos de \textit{Kubernetes} procedentes de cinco fuentes
públicas complementarias, se construyen 599 grafos de clúster originales (1.154
tras aumentación de datos) con aristas que representan relaciones de privilegio,
movimiento lateral, proximidad, espacio de nombres y \acs{RBAC}. La Capa~1
alcanza una F1 macro de 0,9946 en la clasificación de manifiestos individuales, con
\acs{AUC}\,=\,0,9986 y cero falsos negativos en el conjunto de test. La Capa~2,
basada en \acs{GAT} de 3 capas con \textit{embedding} de tipos de arista, logra
F1~macro\,=\,0,803\,\(\pm\)\,0,137
y Precisión@5\,=\,0,720\,\(\pm\)\,0,160 en validación cruzada de 5~pliegues con un
protocolo de particionado consciente de grupos y de familias de plantilla,
superando el objetivo de diseño de P@5\,>\,0,70.

En el conjunto de test independiente —86~grafos con 5~cadenas de ataque de
repositorios reales— el sistema completo, con una
restricción de viabilidad estructural, alcanza Precisión@1\,=\,1,00 y
Precisión@3\,=\,1,00, recupera 4 de las 5~cadenas en las primeras 5~posiciones
(Precisión@5\,=\,0,80) y no produce falsos positivos de clústeres limpios
(\acs{FPR}\_clean\,=\,0\,\%). Una evaluación de usabilidad con tres expertos
otorga una puntuación \acs{SUS} media de 88,3 («excelente»).

El trabajo aporta además una contribución metodológica: la detección y
corrección de una fuga de información por familias de plantilla, junto con la
ablación que la explica. Esta corrección documenta un riesgo relevante para
trabajos futuros con datasets de variantes casi duplicadas. Los resultados
confirman la viabilidad de \textit{kubescan} como control de seguridad
preventivo integrable en \acs{CICD}.

\medskip
\noindent\textbf{Palabras clave:} Kubernetes, Graph Neural Networks, Random Forest,
cadenas de ataque, infraestructura como código.

% ============================================================
\chapter{Abstract}
% ============================================================

This work presents \textit{kubescan}, an automatic attack-chain detection system
for \textit{Kubernetes} infrastructures that combines a Random Forest classifier
(Layer~1) with a \textit{Graph Attention Network} (\acs{GAT}) (Layer~2) and
Genetic Algorithm ensemble weight optimisation (Layer~3).

From 2,827 Kubernetes manifests drawn from five complementary public sources, 599
original cluster graphs are built (1,154 after data augmentation), with edges
representing privilege-escalation, lateral-movement, proximity, namespace, and
\acs{RBAC} relationships. Layer~1 achieves macro-F1\,=\,0.9946, \acs{AUC}\,=\,0.9986,
and zero false negatives on the test set. Layer~2, a 3-layer \acs{GAT} with
edge-type embeddings, achieves
macro-F1\,=\,0.803\,\(\pm\)\,0.137 and Precision@5\,=\,0.720\,\(\pm\)\,0.160 under a
5-fold cross-validation protocol that is group- and template-family-aware,
exceeding the P@5\,>\,0.70 design target.

On the independent test set —
86~graphs containing 5~attack chains from real-world repositories — the full
system, with a structural feasibility constraint,
achieves Precision@1\,=\,1.00 and Precision@3\,=\,1.00, recovers 4 of the
5~attack chains within the top-5 ranking (Precision@5\,=\,0.80), and produces
no false positives on clean clusters (FPR\_clean\,=\,0). A usability
evaluation with three experts yields a mean \acs{SUS} score of 88.3
(``excellent'' rating).

The work additionally contributes a methodological
finding: the detection and correction of template-family information leakage,
together with the ablation that explains it. This
correction documents a risk relevant to future work with datasets of
near-duplicate variants. The results confirm the viability of \textit{kubescan}
as a preventive security control integrable into \acs{CICD}.

\medskip
\noindent\textbf{Keywords:} Kubernetes, Graph Neural Networks, Random Forest,
attack chains, infrastructure as code.
